% ICCIT 2026 draft manuscript - Cyber-Resilient, Explainable Digital Twin
% Framework for Predictive and Cost-Aware Power Quality Management in
% Renewable Microgrids
%
% STATUS: first full-structure draft, generated 2026-08-21 from the
% repo's verified results (README.md Parts 5-6, results/*.csv,
% figures/*.png). NOT compile-tested locally - no LaTeX toolchain was
% available in the environment this was written in. Before submitting:
%   1. Compile this (Overleaf or a local TeX install) and fix any
%      syntax issues - it has not been rendered even once.
%   2. Check ICCIT's actual IEEE Xplore author kit / template and swap
%      \documentclass[conference]{IEEEtran} for theirs if different.
%   3. Fill in real author names/affiliations below.
%   4. Section II (Related Work) cites 7 papers, each verified against
%      Crossref/Dryad metadata (title/authors/journal/DOI all confirmed
%      to resolve) and summarized from structured abstracts, not from a
%      full read of each PDF. Skim the actual papers yourselves before
%      submitting to confirm the framing here is fair to each one.
%   5. The data-provenance sentence in Section III-A is deliberately
%      cautious pending the Jamalpur handoff document your teammate
%      requested (see docs/DATA_PROVENANCE_AND_QUALITY.md) - update it
%      once that lands, in either direction.
%   6. Every number below was pulled from a script that is committed
%      and re-runnable in this repo (see the file cited in each
%      subsection) - if you change a script, regenerate the number.
%   7. Section I now opens with real, cited Bangladesh 2026 load-shedding
%      and Rooppur nuclear plant figures (paper/references.bib,
%      dailystar*/dailysun2026rural/tbs2026rooppur) - verified via web
%      search cross-checked across outlets, not a full article fetch
%      (both outlets 403'd direct fetches). These are news figures that
%      move week to week and Rooppur's commissioning date was still
%      tentative as of the source articles - re-check both right before
%      submission and update the numbers/citations if they've moved.

\documentclass[conference]{IEEEtran}
\usepackage{cite}
\usepackage{amsmath,amssymb}
\usepackage{graphicx}
\usepackage{booktabs}
\usepackage{url}

\begin{document}

\title{A Cyber-Resilient, Explainable Digital Twin Framework for\\
Predictive and Cost-Aware Power Quality Management\\
in Renewable Microgrids}

\author{
\IEEEauthorblockN{Author Name(s) -- TODO}
\IEEEauthorblockA{Department of Electrical \& Electronic Engineering /
Computer Science \& Engineering\\
Institution -- TODO, Bangladesh\\
Email: TODO}
}

\maketitle

\begin{abstract}
Renewable-penetrated microgrids face power quality challenges --
harmonic distortion, voltage sags, and combined weather-electrical
disturbances -- that conventional monitoring addresses only
reactively, without interpretability, cyber-resilience, or economic
context. This paper presents a software-defined digital twin
framework that unifies a physics-based electrical simulation (an
Active Power Filter suppressing current total harmonic distortion
from 27.13\% to 2.07\%) with a machine-learning pipeline trained
entirely on field-measured power-quality data: an unsupervised
Isolation Forest anomaly/cyber-resilience layer, a four-class
disturbance-type predictor (99.7\% held-out accuracy), and a SHAP
explainability layer that confirms the model's decisions align with
known power-engineering causes. Critically, we validate that the two
halves of the twin actually interoperate: waveforms generated by the
physics simulation are classified by the ML model -- trained only on
real data, never on anything simulated -- at 93\% accuracy, after
diagnosing and correcting a zero-source-impedance modeling gap the
test itself exposed. Economic impact is quantified two independent
ways: a physics-based battery-degradation cost model with Monte Carlo
uncertainty bounds (26,955 BDT/yr, 90\% CI [19,623, 34,171]), and a
real HOMER Pro system-level optimization showing renewable-heavy
architectures cost roughly one-quarter as much per kWh as
fossil/battery-only alternatives (\$0.0278 vs.\ \$0.1032/kWh at 81.4\%
vs.\ 0\% renewable fraction). Every result in this paper is produced
by an openly available, re-runnable pipeline, and every dataset's
provenance and known limitations are stated explicitly rather than
implied.
\end{abstract}

\begin{IEEEkeywords}
digital twin, power quality, renewable microgrid, explainable AI,
anomaly detection, active power filter, techno-economic analysis
\end{IEEEkeywords}

\section{Introduction}
% TODO: expand with real citations. Raw material below is adapted from
% the project proposal but needs literature grounding, not just
% restating the proposal's own claims as if they were established facts.

Bangladesh's 2026 power crisis illustrates that this is not simply a
problem of insufficient generation capacity. In early August 2026,
Power Grid Company of Bangladesh data showed the national electricity
shortage climbing from 335 MW to 2{,}167 MW over three days, against a
record shortfall of 2{,}500 MW that pushed large parts of the country
into load-shedding of up to 7--8 hours daily in rural areas and 2--3
hours in urban areas outside Dhaka
\cite{dailystar2026loadshedding,dailysun2026rural}. Yet Bangladesh
Power Development Board data put installed generation capacity at
28{,}919 MW against a projected 2026 summer peak demand of only
18{,}500 MW -- the shortfall traces instead to a gas-supply
constraint (900 MMcf/d available against a stated need of 1{,}200
MMcf/d to run existing plants at peak) \cite{dailystar2026pdbplan}.
Nor is the problem an absence of new capacity in the pipeline: the
country's first nuclear facility, the 1{,}200 MW Rooppur Unit 1,
completed reactor fuel loading in May 2026 and entered final-stage
testing toward commercial operation later in the year
\cite{dailystar2026rooppur,tbs2026rooppur}, yet the record shortfall
above occurred in the same year. Large, centralized capacity
additions -- fossil or nuclear -- do not by themselves solve a crisis
that is substantially about fuel logistics, dispatch, and
distribution reliability, which is exactly the layer at which
decentralized, renewable-heavy microgrids with local storage and
predictive management can help, provided their power quality,
reliability, and cost can be characterized with the same rigor
expected of a conventional grid asset. This paper's contribution
targets precisely that characterization problem, and is framed for
applicability beyond Bangladesh alone (Section~\ref{sec:external-validation}).

The increasing penetration of renewable energy sources and nonlinear
loads in modern microgrids -- the same technology positioned to help
-- introduces its own power quality challenges: voltage instability,
harmonic distortion, and frequency deviation -- that degrade
reliability and accelerate battery wear. Existing
monitoring systems are largely reactive: they report a disturbance
after it has already occurred, offer no explanation for automated
decisions, provide no defense against corrupted or spoofed sensor
readings, and stop short of connecting power-quality degradation to
its economic consequence.

This paper addresses these four gaps with a single integrated
framework combining a physics-based digital twin (electrical
simulation with active harmonic mitigation) and a machine-learning
pipeline (cyber-resilient anomaly detection, disturbance-type
prediction, and explainability), trained and evaluated on
field-measured power-quality data, and cross-validated against the
physics simulation itself -- not just against a held-out split of the
same dataset.

\textbf{Contributions:}
\begin{enumerate}
\item A physics-based power-quality digital twin (weather-driven PV,
  nonlinear-load harmonic injection, an SRF Active Power Filter, and
  an FFT-based THD analyzer) validated against a real annual HOMER
  Pro microgrid simulation, not just a single arbitrary demo timeline.
\item An unsupervised anomaly-detection layer whose validity is
  confirmed by a null-result control: the same method applied to a
  dataset with a known, by-construction-random fault label collapses
  to chance-level ROC-AUC (0.50), while it reaches 0.70 on the real
  disturbance data -- evidence the pipeline is not simply producing a
  plausible-looking number.
\item A four-class disturbance predictor with SHAP explainability,
  whose feature attributions align with power-engineering expectation
  (\texttt{voltage\_rms\_V} dominates Voltage\_Sag predictions).
\item A sim-to-real bridge test: physics-simulated waveforms,
  classified by a model trained only on real field data, at 93\%
  accuracy -- the first evidence in this project that the EEE and CSE
  halves of the "digital twin" actually function as one system rather
  than two independent demonstrations.
\item Two independent, uncertainty-quantified cost analyses (a
  physics-based twin with Monte Carlo bounds, and a real HOMER
  techno-economic optimization), replacing a disputed,
  formula-derived cost column from the raw dataset.
\end{enumerate}

\section{Related Work}

Simulation-based power-quality (PQ) analysis of renewable microgrids
is well established. \cite{hernandezmayoral2024pq} build a detailed
IEEE 14-bus hybrid microgrid in MATLAB/Simulink -- PV, wind, diesel,
storage, and both balanced and unbalanced loads -- and evaluate
voltage/current THD against IEEE-519, concluding that wind generation
and nonlinear loads are dominant harmonic-distortion drivers; this
matches our own finding that the 6-pulse rectifier load is the sole
harmonic source worth mitigating in our twin. Compensation and control
have increasingly moved toward AI-tuned rather than fixed-parameter
designs: \cite{john2026aipqreview} review over 100 studies applying
swarm intelligence and hybrid AI (ANFIS, PSO-ANN, adaptively-tuned
APF/D-STATCOM/UPQC) for harmonic mitigation, reporting faster
convergence and lower THD than fixed-parameter compensators, and
\cite{prakash2025aidcmicrogrid} pair an ANN radial-basis MPPT
controller with dual fractional-order PID for DC-microgrid battery
management. Our twin's SRF-based active power filter (Section
\ref{sec:twin}) is, by this standard, a fixed-parameter design;
adaptively tuning it is identified as future work in Section
\ref{sec:limitations} rather than claimed here.

On the forecasting/classification side, \cite{jahan2026mlpq} benchmark
nine ML/DL models (LSTM, GRU, DNN, CNN1D-LSTM, BiLSTM, attention-RNN,
decision tree, SVM, XGBoost) predicting PQ parameters -- voltage,
THDu, THDi, short-term flicker -- on a real off-grid microgrid
platform, finding deep models more accurate and classical ones far
cheaper to run. That work forecasts continuous PQ metrics; ours
instead classifies discrete disturbance \emph{type} (Section
\ref{sec:twin}) and adds SHAP-based explainability plus a
cyber-resilience layer, a different but complementary problem framing.
\cite{osifeko2025aiforecasting} benchmark 12 models on five years of
South African grid data and explicitly document the constraints of
forecasting in data-sparse, developing-region deployments -- limited
local training data, missing real-time measurement coverage, and poor
transfer of models trained on developed-country data. Bangladesh's
Jamalpur dataset used here faces the same category of constraint,
which is part of why this paper reports every accuracy number
alongside an explicit uncertainty or robustness check (bootstrap CIs,
a null-result control, sim-to-real transfer) rather than a bare point
estimate.

On techno-economic optimization, \cite{daniel2025dualopt} propose
ALA-TKAN, a single AI framework unifying renewable/load forecasting
with dual-optimization energy management, reporting THD of 1.4\% and
energy cost of \$0.8/Wh for a hybrid PV-wind-fuel-cell-battery
microgrid. Section~\ref{sec:cost} takes a deliberately different
approach: rather than one model producing both a forecast and a cost
figure, we keep the physics twin's Monte Carlo cost projection and
HOMER Pro's independent commercial optimizer separate and both
interval-reported, specifically so an assumed-parameter model is never
mistaken for a measured or independently-optimized one -- a provenance
concern raised throughout this paper (Section~\ref{sec:limitations}).

Beyond the datasets used directly in this work (Section III-A),
real, published microgrid instrumentation datasets exist elsewhere in
the literature; for example, \cite{bashir2023mesadelsol} report
power, voltage, frequency, and thermal telemetry from the Mesa Del
Sol microgrid (Albuquerque, USA) at 10-second resolution over 15
months. That dataset does not include harmonic/THD measurements, so
it cannot substitute for the Jamalpur data used here or validate
this paper's disturbance classifier, but we use one representative
month of it (April 2023, 259{,}200 valid samples) as an independent,
unrelated real installation against which to sanity-check the
digital twin's voltage/frequency assumptions -- see
Section~\ref{sec:external-validation}.

None of the above combine a physics-based digital twin, a classifier
trained exclusively on real field data and stress-tested against
simulated waveforms, cyber-resilient anomaly detection validated with
a null-result control, and explicit cost uncertainty quantification
into a single closed-loop system; that combination is this paper's
contribution relative to the reviewed literature.

\section{Methodology}

\subsection{Data}
Two datasets support this work. The first is a 5,000-row,
single-phase power-quality dataset (voltage, current, frequency,
harmonics, THD, weather, and a sensor-fault flag), reported by a
project team member to originate from field measurements at an
industrial site in Jamalpur, Bangladesh; the original collection
protocol was not independently available to the authors at the time
of writing, and independent confirmation is in progress. This
caveat is stated because two derived columns in the same dataset
(\texttt{economic\_cost\_BDT}, \texttt{battery\_capacity\_loss\_pct})
were found to be exact algebraic functions of a third column rather
than independent measurements, and are therefore excluded from this
paper's cost analysis (Section~\ref{sec:cost}) as a matter of
methodological hygiene, independent of the provenance question.

The second is a real HOMER Pro annual hourly simulation (8,760 rows:
solar, wind, load, battery state-of-charge, and system economics)
together with HOMER's own multi-architecture optimization output,
used to drive realistic operating conditions for the physics twin
(Section~\ref{sec:twin}) and to provide an independent,
provenance-clean cost comparison (Section~\ref{sec:cost}).

A third, 50,000-row synthetic dataset generated for pipeline
development is explicitly \emph{not} used to support any claim in
this paper; it is watermarked at generation time and its own
repository documentation states it must not be presented as
real-world evidence.

\subsection{Physics-Based Digital Twin}
\label{sec:twin}
The electrical-level twin models a Bangladesh-context PV microgrid:
irradiance- and NOCT-derated-temperature-driven PV output, a
nonlinear six-pulse-rectifier-type load injecting the characteristic
non-triplen harmonic set (5th, 7th, 11th, 13th), a Synchronous
Reference Frame (Park-transform) Active Power Filter, and an
FFT-based THD analyzer over 5-cycle (100\,ms) windows at a 10\,kHz
solver rate. The filter suppresses source-current THD from 27.13\%
to 2.07\% within one cycle of activation, comfortably inside the IEEE
519 5\% limit.

To move beyond a single arbitrary demonstration timeline, three
representative operating points were drawn directly from the HOMER
annual simulation: the peak-load hour, the hour with the sharpest
hour-over-hour drop in renewable output, and a low-solar/high-wind/
high-load hour. The Active Power Filter suppresses THD to the same
level at all three points -- an expected consequence of the harmonic
model's fixed percentage-of-fundamental structure, and reported here
as evidence of suppression \emph{robustness} across the annual
operating envelope, not as three scenarios of differing severity.
Projected battery cost, by contrast, varies by an order of magnitude
across the same three hours, driven by each hour's solar throughput.

\subsection{Cyber-Resilient Anomaly Detection}
An Isolation Forest, trained without access to the ground-truth
sensor-fault label, is evaluated on both the real disturbance dataset
and, as a validity control, on the synthetic dataset's fault flag --
which is generated purely at random by construction and therefore
has a true chance-level relationship to every feature. The method
reaches ROC-AUC 0.70 on the real data and collapses to 0.50 (chance)
on the null-control dataset, evidence the pipeline is measuring a
real signal rather than an artifact of the method. A supervised
Random Forest given direct label access reaches ROC-AUC 0.83,
indicating the fault signal is present in the electrical features but
is not shaped like a natural multivariate outlier -- a methodological
finding in its own right, motivating semi-supervised approaches as
future work rather than indicating the unsupervised method "failed."

\subsection{Disturbance Prediction and Explainability}
A four-class Random Forest (None / Voltage\_Sag /
Harmonic\_Distortion / Combined\_Weather\_Electrical) trained on ten
electrical and weather features reaches 99.7\% held-out accuracy
(99.26\% mean over five seeds). SHAP analysis shows
\texttt{voltage\_rms\_V} dominates Voltage\_Sag predictions and
\texttt{THD\_voltage\_pct}/harmonic-order features dominate
Harmonic\_Distortion predictions -- consistent with the physical
mechanism of each disturbance type, not a black-box coincidence.

\subsection{Closing the Loop: Sim-to-Real Validation}
\label{sec:bridge}
The physics twin (Section~\ref{sec:twin}) and the disturbance
classifier (above) were, until this experiment, two independent
pipelines running on unrelated data. We test whether they actually
function as one system by injecting the same four disturbance
classes into the electrical twin, extracting the identical
ten-feature vector, and classifying the \emph{simulated} waveform
with the model trained \emph{only} on real Jamalpur data.

An initial attempt reached 25\% accuracy -- the classifier defaulted
to the majority class for every input. Diagnosis traced this to the
twin's ideal, zero-impedance voltage source: \texttt{THD\_voltage\_pct}
was structurally always exactly zero, yet this is the single largest
discriminator between Harmonic\_Distortion (real-data mean 11.3\%)
and Voltage\_Sag (2.9\%) in the training distribution. We added a
standard linearly-frequency-scaling inductive source reactance
($X_h = h \cdot X_1$) and recalibrated the injected harmonic
amplitude to real-world THD magnitudes rather than an uncalibrated
textbook default (which, at 27\%, exceeded every value in the real
training distribution). After this correction -- itself a
textbook-justified addition, not a curve fit to the test set --
accuracy reached \textbf{93\%} (Table~\ref{tab:bridge}), and the
anomaly detector flagged 96--100\% of simulated disturbed cases
versus 28\% of simulated normal cases.

\begin{table}[t]
\centering
\caption{Sim-to-real bridge: classifier trained only on real data,
evaluated on physics-simulated waveforms (25 repetitions/class)}
\label{tab:bridge}
\begin{tabular}{lccc}
\toprule
Class & Precision & Recall & F1 \\
\midrule
None & 1.00 & 1.00 & 1.00 \\
Harmonic\_Distortion & 1.00 & 1.00 & 1.00 \\
Voltage\_Sag & 0.78 & 1.00 & 0.88 \\
Combined\_Weather\_Electrical & 1.00 & 0.72 & 0.84 \\
\midrule
\textbf{Overall accuracy} & \multicolumn{3}{c}{\textbf{0.93}} \\
\bottomrule
\end{tabular}
\end{table}

\subsection{Cost Quantification}
\label{sec:cost}
Two independent, provenance-clean cost analyses are used in place of
the dataset's disputed \texttt{economic\_cost\_BDT} column
(Section III-A). First, the physics twin's own
battery-degradation-to-cost model, evaluated under a 200-run Monte
Carlo with its two assumed parameters (battery replacement cost,
degradation rate) perturbed $\pm$20\%, gives 26,955 BDT/yr (point
estimate) with a 90\% interval of [19,623, 34,171] BDT/yr -- reported
as an interval rather than a false-precision point estimate. Second,
HOMER Pro's own multi-architecture optimization
(Table~\ref{tab:homer}) shows renewable fraction and cost of energy
are strongly inversely related in this system: an 81.4\%-renewable
configuration reaches \$0.0278/kWh, roughly one-quarter the cost of a
0\%-renewable, fuel-cell-only configuration at \$0.1032/kWh.

\begin{table}[t]
\centering
\caption{HOMER optimization: three candidate architectures}
\label{tab:homer}
\begin{tabular}{lccc}
\toprule
Config & Renewable Frac. & COE (\$/kWh) & NPC (\$) \\
\midrule
1 (PV+WT+FC+BESS) & 81.4\% & 0.0278 & 501,600 \\
2 (WT+FC+BESS)     & 53.3\% & 0.0671 & 736,617 \\
3 (FC+BESS only)   & 0.0\%  & 0.1032 & 1,120,089 \\
\bottomrule
\end{tabular}
\end{table}

\subsection{External Validation Against an Independent Real Microgrid}
\label{sec:external-validation}
The physics twin's \texttt{plant\_step()} models grid voltage and
frequency as exact constants except during one scripted sag event --
an idealization never previously checked against real telemetry. To
test it independently of the Jamalpur dataset and of this project's
own team, we compare it against one representative month (April 2023,
259{,}200 valid 10-second samples, after removing an 8{,}424-row
sentinel/placeholder block inserted by the source logger) of the Mesa
Del Sol dataset \cite{bashir2023mesadelsol} -- an unrelated 60 Hz,
$\approx$484 V US commercial microgrid, so only normalized (per-unit)
behavior, not absolute magnitudes, is comparable to the 50 Hz, 230 V
Jamalpur context. Even with no fault or sag scripted for nearly the
entire month, real frequency shows a continuous standard deviation of
$\approx$0.03\% p.u. (full range $\approx\pm$0.4\% p.u.) and real
voltage shows a continuous standard deviation of $\approx$0.5--0.6\%
p.u. (full range $\approx$4--5\% p.u.) at two independent measurement
points each. This is concrete, independently-sourced evidence that
the twin's constant-V/constant-f-between-events assumption is an
idealization worth stating explicitly (Section~\ref{sec:limitations})
rather than an unexamined convenience. This dataset has no
THD/harmonic columns, so it validates voltage/frequency behavior only
and cannot validate the disturbance classifier or THD results.

\section{Results Summary}
Table~\ref{tab:summary} consolidates the headline results referenced
above; full per-seed statistics, confusion matrices, and SHAP plots
are in the project repository's \texttt{results/} and
\texttt{figures/} directories, generated by the scripts cited in each
row.

\begin{table}[t]
\centering
\caption{Headline results (script that generates each number)}
\label{tab:summary}
\begin{tabular}{p{4.6cm}p{3cm}}
\toprule
Result & Source script \\
\midrule
THD$_i$: 27.13\% $\to$ 2.07\% & \texttt{microgrid\_pq\_twin.py} \\
Disturbance classifier: 99.7\% acc. & \texttt{disturbance\_classifier.py} \\
Anomaly detector: ROC-AUC 0.70 vs.\ 0.50 (null control) & \texttt{anomaly\_detection.py} \\
Sim-to-real bridge: 93\% acc. & \texttt{twin\_to\_cse\_bridge.py} \\
Cost: 26,955 BDT/yr, 90\% CI [19.6k, 34.2k] & \texttt{sensitivity\_analysis.py} \\
HOMER: 81.4\% RE $\to$ \$0.0278/kWh & \texttt{homer\_economic\_optimization.py} \\
External V/f check: real std $\approx$0.03\% (f), $\approx$0.5--0.6\% (V) p.u.\ & \texttt{mesa\_del\_sol\_validation.py} \\
\bottomrule
\end{tabular}
\end{table}

\section{Discussion and Limitations}
\label{sec:limitations}
Several limitations are stated explicitly rather than left implicit.
The disturbance dataset's field-measurement provenance rests on a
team member's confirmation, not an independently verifiable handoff
record at the time of writing; a documentary request is in progress
and this section will be updated once resolved. The near-perfect
disturbance-classification accuracy (99.7\%) is plausible given the
data's exceptionally low noise and complete sensor coverage, both
atypical of field telemetry, and is reported alongside -- not instead
of -- the harder unsupervised anomaly-detection result (F1 $\approx$
0.10) precisely because a reviewer should be able to weigh both. The
93\% sim-to-real bridge result is, in this light, meaningful
independent evidence: a model cannot be overfit to idiosyncrasies of
one dataset it has never seen, so its ability to generalize to an
unrelated first-principles simulation is a genuine (if imperfect --
Voltage\_Sag and Combined\_Weather\_Electrical still show some
cross-confusion) signal that the classifier learned real
power-quality structure. The physics twin itself remains an
idealized, single-branch electrical model (no source impedance in
its primary form, uncited assumed cost parameters, Simulink
cross-validation still pending as of this writing) that holds voltage
and frequency at an exact constant except during one scripted event;
Section~\ref{sec:external-validation}'s comparison against an
independent real microgrid confirms this last point is a real
idealization, not a negligible one -- real installations show small
but continuous V/f variability even under normal operation. Each of
these is addressed either with an explicit uncertainty bound
(Section~\ref{sec:cost}) or a stated, unfixed limitation rather than
silently absorbed into a headline number.

Returning to the motivation in Section~I: nothing in this framework's
methodology is intrinsically Bangladesh-specific. The physics twin's
V/f analysis is normalized (per-unit), which is precisely what let
Section~\ref{sec:external-validation} validate it against an
unrelated 60 Hz, $\approx$484 V US microgrid rather than only against
Jamalpur's 50 Hz, 230 V context; the cost-quantification methodology
(Section~\ref{sec:cost}) substitutes local currency and tariff inputs
for BDT without structural change. A decentralized, digital-twin-driven
approach of this kind is therefore a candidate response on two fronts
at once: for grids like Bangladesh's, where the 2026 data cited in
Section~I show the crisis is substantially about fuel logistics and
dispatch reliability rather than installed capacity
\cite{dailystar2026pdbplan,dailystar2026rooppur}, decentralized
renewable microgrids reduce dependence on that centralized,
fuel-constrained supply chain; and for renewable-heavy grids in
developed contexts, the same interpretable, cost-aware,
cyber-resilient monitoring layer addresses the harmonic and
disturbance-management burden that renewable penetration itself
introduces (Section~II). This paper does not claim to have deployed
or field-tested either case -- only that the underlying methodology
was built and validated in a way that does not foreclose either one.

\section{Conclusion}
This paper presented a digital twin framework for renewable-microgrid
power quality management that is cyber-resilient (validated anomaly
detection with a null-result control), explainable (SHAP attributions
matching power-engineering expectation), and cost-aware (two
independent, uncertainty-quantified cost analyses). Unlike a
side-by-side pairing of a physics simulation and a machine-learning
pipeline, the two halves were shown to function as one system: a
classifier trained exclusively on real field data correctly
identifies disturbance types in physics-simulated waveforms at 93\%
accuracy. Future work includes completing the pending Simulink
cross-validation, extending the source-impedance correction
(Section~\ref{sec:bridge}) into the primary digital twin rather than
the bridge test alone, and revisiting the dataset provenance question
once documentary confirmation is available.

\bibliographystyle{IEEEtran}
\bibliography{references}
% references.bib has ONE verified entry (Mesa Del Sol dataset). Add
% real citations for the rest of Related Work before submitting -
% every entry so far was independently checked, not guessed; keep it
% that way.

\end{document}
